\documentclass[12pt]{article}
\usepackage[T1]{fontenc}
\usepackage[utf8]{inputenc}
\usepackage{lmodern}
\usepackage{textcomp}
\usepackage{enumitem}
\usepackage{geometry}
\usepackage{graphicx}
\usepackage{amsmath, amssymb}
\geometry{margin=1in}
\begin{document}
\begin{center}
Astronomy - Graduate Level Assessment 18 \
Total Marks: 40 \
Answer all questions.
\end{center}

\section*{Multiple Choice (10 marks)}
\begin{enumerate}[label=\arabic*.]
\item Which of the following is/are NOT caused by orbital resonance?
\begin{enumerate}[label=(\alph*)]
    \item 2:3 periodic ratio of Neptune:Pluto
    \item Kirkwood Gaps.
    \item Gaps in Saturn's rings.
    \item Breaking of small Jovian moons to form ring materials.
\end{enumerate}
\item Why can't we see the Sun's corona except during total solar eclipses?
\begin{enumerate}[label=(\alph*)]
    \item The corona is made up mostly of charged particles which aren't luminous.
    \item It's much too cool to emit visible light
    \item We can't see magnetic fields
    \item It's too diffuse
\end{enumerate}
\item Meteorites with high metal content probably are
\begin{enumerate}[label=(\alph*)]
    \item pieces of comets rather than of asteroids.
    \item chunks of large differentiated asteroids that were shattered by collisions.
    \item chunks of rock chipped off the planet Mars.
    \item leftover chunks of rock from the earliest period in the formation of the solar system.
\end{enumerate}
\item We were first able to accurately measure the diameter of Pluto from:
\begin{enumerate}[label=(\alph*)]
    \item a New Horizons flyby in the 1990 s
    \item Hubble Space Telescope images that resolved Pluto's disk
    \item brightness measurements made during mutual eclipses of Pluto and Charon
    \item radar observations made by the Arecibo telescope
\end{enumerate}
\item Which of the following characteristics would not necessarily suggest that a rock we found is a meteorite.
\begin{enumerate}[label=(\alph*)]
    \item It has a fusion crust
    \item It contains solidified spherical droplets
    \item It is highly processed
    \item It has different elemental composition than earth
\end{enumerate}
\end{enumerate}

\section*{True / False (10 marks)}
\textit{Answer True or False}
\begin{enumerate}[label=\arabic*.]
\setcounter{enumi}{5}
\item True or False: Five kilograms of feathers would weigh the most on the moon due to its greater mass compared to other options.
\item True or False: Knowing both the actual brightness and apparent brightness of an object allows for the estimation of its distance from the observer.
\item True or False: The axis of the Earth is tilted at an angle of approximately 23.5 degrees relative to its orbital plane around the Sun.
\item True or False: Centaurus is often mistaken for a W-shaped constellation, but it is actually more complex in its structure and appearance.
\item True or False: Dark energy is a theoretical model that explains the observed acceleration of the universe's expansion.
\end{enumerate}

\section*{Long Form Response (20 marks)}
\textit{Answer each question with a clear and complete explanation.}
\begin{enumerate}[label=\arabic*.]
\setcounter{enumi}{10}
\item Discuss the changes in the emission spectrum of a blackbody when its temperature is reduced to half of its original value, focusing on power emitted and peak emission wavelength.
\item Discuss the significance of constellations in relation to the Milky Way. In your analysis, identify at least one constellation that is not located along the Milky Way and explain its relevance.
\end{enumerate}

\vfill
\noindent\textit{End of Paper}
\end{document}